% Options for packages loaded elsewhere
\PassOptionsToPackage{unicode}{hyperref}
\PassOptionsToPackage{hyphens}{url}
%
\documentclass[
  ignorenonframetext,
]{beamer}
\usepackage{pgfpages}
\setbeamertemplate{caption}[numbered]
\setbeamertemplate{caption label separator}{: }
\setbeamercolor{caption name}{fg=normal text.fg}
\beamertemplatenavigationsymbolsempty
% Prevent slide breaks in the middle of a paragraph
\widowpenalties 1 10000
\raggedbottom
\setbeamertemplate{part page}{
  \centering
  \begin{beamercolorbox}[sep=16pt,center]{part title}
    \usebeamerfont{part title}\insertpart\par
  \end{beamercolorbox}
}
\setbeamertemplate{section page}{
  \centering
  \begin{beamercolorbox}[sep=12pt,center]{part title}
    \usebeamerfont{section title}\insertsection\par
  \end{beamercolorbox}
}
\setbeamertemplate{subsection page}{
  \centering
  \begin{beamercolorbox}[sep=8pt,center]{part title}
    \usebeamerfont{subsection title}\insertsubsection\par
  \end{beamercolorbox}
}
\AtBeginPart{
  \frame{\partpage}
}
\AtBeginSection{
  \ifbibliography
  \else
    \frame{\sectionpage}
  \fi
}
\AtBeginSubsection{
  \frame{\subsectionpage}
}

\usepackage{amsmath,amssymb}
\usepackage{iftex}
\ifPDFTeX
  \usepackage[T1]{fontenc}
  \usepackage[utf8]{inputenc}
  \usepackage{textcomp} % provide euro and other symbols
\else % if luatex or xetex
  \usepackage{unicode-math}
  \defaultfontfeatures{Scale=MatchLowercase}
  \defaultfontfeatures[\rmfamily]{Ligatures=TeX,Scale=1}
\fi
\usepackage{lmodern}
\ifPDFTeX\else  
    % xetex/luatex font selection
\fi
% Use upquote if available, for straight quotes in verbatim environments
\IfFileExists{upquote.sty}{\usepackage{upquote}}{}
\IfFileExists{microtype.sty}{% use microtype if available
  \usepackage[]{microtype}
  \UseMicrotypeSet[protrusion]{basicmath} % disable protrusion for tt fonts
}{}
\makeatletter
\@ifundefined{KOMAClassName}{% if non-KOMA class
  \IfFileExists{parskip.sty}{%
    \usepackage{parskip}
  }{% else
    \setlength{\parindent}{0pt}
    \setlength{\parskip}{6pt plus 2pt minus 1pt}}
}{% if KOMA class
  \KOMAoptions{parskip=half}}
\makeatother
\usepackage{xcolor}
\newif\ifbibliography
\setlength{\emergencystretch}{3em} % prevent overfull lines
\setcounter{secnumdepth}{-\maxdimen} % remove section numbering


\providecommand{\tightlist}{%
  \setlength{\itemsep}{0pt}\setlength{\parskip}{0pt}}\usepackage{longtable,booktabs,array}
\usepackage{calc} % for calculating minipage widths
\usepackage{caption}
% Make caption package work with longtable
\makeatletter
\def\fnum@table{\tablename~\thetable}
\makeatother
\usepackage{graphicx}
\makeatletter
\def\maxwidth{\ifdim\Gin@nat@width>\linewidth\linewidth\else\Gin@nat@width\fi}
\def\maxheight{\ifdim\Gin@nat@height>\textheight\textheight\else\Gin@nat@height\fi}
\makeatother
% Scale images if necessary, so that they will not overflow the page
% margins by default, and it is still possible to overwrite the defaults
% using explicit options in \includegraphics[width, height, ...]{}
\setkeys{Gin}{width=\maxwidth,height=\maxheight,keepaspectratio}
% Set default figure placement to htbp
\makeatletter
\def\fps@figure{htbp}
\makeatother

\makeatletter
\@ifpackageloaded{caption}{}{\usepackage{caption}}
\AtBeginDocument{%
\ifdefined\contentsname
  \renewcommand*\contentsname{Table of contents}
\else
  \newcommand\contentsname{Table of contents}
\fi
\ifdefined\listfigurename
  \renewcommand*\listfigurename{List of Figures}
\else
  \newcommand\listfigurename{List of Figures}
\fi
\ifdefined\listtablename
  \renewcommand*\listtablename{List of Tables}
\else
  \newcommand\listtablename{List of Tables}
\fi
\ifdefined\figurename
  \renewcommand*\figurename{Figure}
\else
  \newcommand\figurename{Figure}
\fi
\ifdefined\tablename
  \renewcommand*\tablename{Table}
\else
  \newcommand\tablename{Table}
\fi
}
\@ifpackageloaded{float}{}{\usepackage{float}}
\floatstyle{ruled}
\@ifundefined{c@chapter}{\newfloat{codelisting}{h}{lop}}{\newfloat{codelisting}{h}{lop}[chapter]}
\floatname{codelisting}{Listing}
\newcommand*\listoflistings{\listof{codelisting}{List of Listings}}
\makeatother
\makeatletter
\makeatother
\makeatletter
\@ifpackageloaded{caption}{}{\usepackage{caption}}
\@ifpackageloaded{subcaption}{}{\usepackage{subcaption}}
\makeatother

\ifLuaTeX
  \usepackage{selnolig}  % disable illegal ligatures
\fi
\usepackage{bookmark}

\IfFileExists{xurl.sty}{\usepackage{xurl}}{} % add URL line breaks if available
\urlstyle{same} % disable monospaced font for URLs
\hypersetup{
  pdftitle={Parameter estimation with the Lunar Gravitational Wave Antenna},
  pdfauthor={Jacopo Tissino},
  hidelinks,
  pdfcreator={LaTeX via pandoc}}


\title{Parameter estimation with the Lunar Gravitational Wave Antenna}
\subtitle{Moriond 2025}
\author{Jacopo Tissino}
\date{2025-03-31}

\begin{document}
\frame{\titlepage}


\begin{frame}{What is LGWA?}
\phantomsection\label{what-is-lgwa}
\begin{center}
\includegraphics[width=\textwidth,height=0.75\textheight]{../figures/LGWA_crater_lighter.png}
\end{center}

\footnotesize Figure from the LGWA collaboration.
\end{frame}

\begin{frame}{LGWA payload}
\phantomsection\label{lgwa-payload}
\begin{center}
\includegraphics[width=\textwidth,height=0.75\textheight]{../figures/LGWApayload.png}
\end{center}

\footnotesize Figure from
\href{https://pubs.aip.org/aip/jap/article/133/24/244501/2899601/The-payload-of-the-Lunar-Gravitational-wave}{van
Heijningen et al., 2023}.
\end{frame}

\begin{frame}{LGWA readout}
\phantomsection\label{lgwa-readout}
\begin{center}
\includegraphics[width=\textwidth,height=0.75\textheight]{../figures/squid-readout.png}
\end{center}

\footnotesize Figure from
\href{https://pubs.aip.org/aip/jap/article/133/24/244501/2899601/The-payload-of-the-Lunar-Gravitational-wave}{van
Heijningen et al., 2023}.
\end{frame}

\begin{frame}{Gravitational wave response}
\phantomsection\label{gravitational-wave-response}
How does this readout noise translate to noise on our measurement of
gravitational waves?

We assume that the response function from the gravitational
polarizations \(h_{ij}(t)\) to the displacement at the sensor \(s(t)\)
is \emph{linear}: in the frequency domain,

\[ s(f) = L(f) D_{ij} h_{ij}(f)
\]

\begin{itemize}
\tightlist
\item
  \(L(f)\) is the \textbf{response function}: with the dimensions of a
  length, is the displacement amplitude corresponding to a monochromatic
  unit strain;
\item
  \(D_{ij} = n_i b_j\) is a \textbf{pattern function}, where \(b_j\) is
  a tangent unit vector to the surface while \(n_i\) is normal to it.
\end{itemize}
\end{frame}

\begin{frame}{Structure of LGWA}
\phantomsection\label{structure-of-lgwa}
LGWA is planned to have:

\begin{itemize}
\tightlist
\item
  4 \textbf{stations}, to reduce overall noise and perform active
  cancellation;
\item
  for each station, two orthogonal \textbf{horizontal measurement
  channels}.
\end{itemize}
\end{frame}

\begin{frame}{Strain noise}
\phantomsection\label{strain-noise}
If we have:

\begin{itemize}
\tightlist
\item
  an estimate for the PSD of the detector \textbf{displacement noise},
  \(S_\text{disp}(f)\),
\item
  an estimate of the lunar \textbf{response} \(L(f)\),
\item
  the number of stations \(N\),
\end{itemize}

we can compute the strain noise as

\[ S_{\text{strain}} (f) = \frac{S _{\text{disp}}(f)}{N L^2(f)}
\]
\end{frame}

\begin{frame}{LGWA estimated noise}
\phantomsection\label{lgwa-estimated-noise}
\begin{center}
\includegraphics[width=\textwidth,height=0.75\textheight]{../figures/LGWA_sensitivity.png}
\end{center}
\end{frame}

\begin{frame}{The multiband concept}
\phantomsection\label{the-multiband-concept}
\begin{center}
\includegraphics[width=\textwidth,height=0.75\textheight]{../figures/Multiband_concept.pdf}
\end{center}
\end{frame}

\begin{frame}{The reference frame problem}
\phantomsection\label{the-reference-frame-problem}
The frequency-domain waveform at the detector has a travel time factor
from the center of the reference frame:

\[ h(f) \propto \exp( i \vec{k} \cdot (\vec{r} - \vec{r}_0))
\]

The center \(\vec{r}_0\) should be rigidly connected with the SSB, but
it can be elsewhere!
\end{frame}

\begin{frame}{Sky position measurement - 1 month of observation}
\phantomsection\label{sky-position-measurement---1-month-of-observation}
\includegraphics{lgwa-pe-presentation-moriond_files/figure-beamer/cell-4-output-1.pdf}
\end{frame}

\begin{frame}{Sky position measurement - 1 year}
\phantomsection\label{sky-position-measurement---1-year}
\includegraphics{lgwa-pe-presentation-moriond_files/figure-beamer/cell-5-output-1.pdf}
\end{frame}

\begin{frame}{Sky position measurement - early warning}
\phantomsection\label{sky-position-measurement---early-warning}
\includegraphics{lgwa-pe-presentation-moriond_files/figure-beamer/cell-6-output-1.pdf}
\end{frame}

\begin{frame}{Mass constraints}
\phantomsection\label{mass-constraints}
\includegraphics{lgwa-pe-presentation-moriond_files/figure-beamer/cell-7-output-1.pdf}
\end{frame}

\begin{frame}{Sky localization}
\phantomsection\label{sky-localization}
\includegraphics{lgwa-pe-presentation-moriond_files/figure-beamer/cell-8-output-1.pdf}
\end{frame}

\begin{frame}{Phase constraints}
\phantomsection\label{phase-constraints}
\includegraphics{lgwa-pe-presentation-moriond_files/figure-beamer/cell-10-output-1.pdf}
\end{frame}

\begin{frame}{Open problems: noise and moonquakes}
\phantomsection\label{open-problems-noise-and-moonquakes}
\begin{itemize}
\tightlist
\item
  Can we demonstrate active noise cancellation?
\item
  LGWA data will have gaps when moonquakes get too strong: how much do
  they impact PE?
\end{itemize}
\end{frame}

\begin{frame}{Open problems: response calibration}
\phantomsection\label{open-problems-response-calibration}
\begin{itemize}
\tightlist
\item
  We are assuming a known lunar response.
\item
  How does uncertainty there translate into parameter estimation?
\item
  Can we demonstrate the ``multibanding calibration trick''?
\end{itemize}
\end{frame}

\begin{frame}{Open problems: parameter estimation}
\phantomsection\label{open-problems-parameter-estimation}
\begin{itemize}
\tightlist
\item
  likelihood acceleration (relative binning or similar)
\item
  analytical marginalizations (e.g.~in phase)
\item
  different sources

  \begin{itemize}
  \tightlist
  \item
    IMBH
  \item
    stellar BBH
  \item
    DWD
  \item
    \ldots{}
  \end{itemize}
\end{itemize}
\end{frame}




\end{document}
